% GENERATED_BY: oc_core_monograph_translation_draft_pipeline_v1.py
% AI_ASSISTED_TRANSLATION_DRAFT: TRUE
% THEOREM_REVIEW_STATUS: PENDING
% THEOREM_REVIEW_MODE: FORMAL_AI_GATE__CODEX_CLI_CHATGPT
% THEOREM_REVIEW_REVIEWER_ID:
% THEOREM_REVIEW_AUTHORITY_CLASS:
% THEOREM_REVIEWED_AT:
% THEOREM_REVIEW_DOSSIER_REF:
% SOURCE_LANGUAGE: EN
% TARGET_LANGUAGE: RU
% SOURCE_EN_REF: content/m_spaces/m6.tex
% ================================================================
% ==== FILE: content/m_spaces/m6.tex
% ================================================================

% ==============================
%  Ontology of Continua — Core
%  Meta-Space: M6
%  FULL MODULE — FINAL VERSION
% ==============================

\subsubsection{\texorpdfstring{$M_6$}{M_6} Обзор}
\label{sec:m6-overview}

$M_6$ — это мета-пространство, которое делает \emph{когнитивную организацию} допустимой.
Если $M_5$ поддерживает электрическую возбудимость и спайк-циклы, то $M_6$ обеспечивает:

\begin{itemize}
    \item устойчивые аттракторы над последовательностями спайков,
    \item семантические оси и геометрию репрезентации,
    \item протологические операции, действующие на нейронные состояния,
    \item динамику классификации и прогнозирования,
    \item возникновение символических и подсемантических понятий,
    \item ограничения когерентности, необходимые для когнитивных континуумов ($K_6$).
\end{itemize}

Таким образом, $M_6$ — минимальная среда для \emph{смыслонасыщенной нейронной динамики}.

% ---------------------------------------------------------------
\subsubsection*{1. Допустимое пространство конфигураций \texorpdfstring{$\Omega(M_6)$}{\Omega(M_6)}}

\[
\Omega(M_6) =
\left\{
(S(t),\ A_{\mathrm{sem}},\
\mathcal{W},\
\mathcal{A},\
C_{\mathrm{att}},\
\Pi_6)
\ \big|\ \text{выполняются условия допустимости}
\right\}.
\]

Здесь:

\begin{itemize}
    \item $S(t)$ — траектории нейронных состояний (последовательности спайков, частоты разрядов, фазовые коды),
    \item $\mathcal{W}$ — тензор синаптических весов с ограничениями пластичности,
    \item $\mathcal{A}$ — допустимое множество аттракторов,
    \item $A_{\mathrm{sem}}$ — семантические оси, встроенные в геометрию аттрактора,
    \item $C_{\mathrm{att}}$ — семейство циклов внимания / рабочей памяти,
    \item $\Pi_6$ — ограничения, обеспечивающие устойчивость и разделимость репрезентаций.
\end{itemize}

Конфигурация допустима, если:

\begin{enumerate}
    \item нейронная динамика может сходиться хотя бы к одному устойчивому аттрактору;
    \item бассейны притяжения разделимы гиперплоскостями или многообразиями;
    \item отображение между входными паттернами и аттракторами согласовано;
    \item правила пластичности сохраняют когерентность, а не разрушают структуры.
\end{enumerate}

% ---------------------------------------------------------------
\subsubsection*{2. Допустимые оси \texorpdfstring{$A(M_6)$}{A(M_6)}}

$M_6$ вводит новый класс осей:

\[
A_{\mathrm{concept}} = \{ \text{концептуальные различия, возникающие из аттракторов} \}.
\]

Эти оси соответствуют \emph{смыслонасыщенным степеням свободы}.

Другие допустимые оси:

\begin{itemize}
    \item $A_{\mathrm{sem}}$ — оси семантического вложения в нейронном пространстве,
    \item $A_{\mathrm{wm}}$ — ось поддержания рабочей памяти,
    \item $A_{\mathrm{att}}$ — ось модуляции внимания,
    \item $A_{\mathrm{proto\_logic}}$ — оси, поддерживающие протологические операции,
    \item $A_{\mathrm{prediction}}$ — оси предиктивного кодирования.
\end{itemize}

Необходимое условие для $K_6$:

\[
\dim(A_{\mathrm{sem}}) \ge 1
\quad\text{и}\quad
A_{\mathrm{concept}} \subseteq A(M_6).
\]

% ---------------------------------------------------------------
\subsubsection*{3. Потенциалы \texorpdfstring{$P(M_6)$}{P(M_6)}}

Когнитивные потенциалы обобщают электро-динамические потенциалы $M_5$ на более высокий уровень структуры:

\[
P(M_6) =
(U_{\mathrm{att}},\
U_{\mathrm{wm}},\
U_{\mathrm{concept}},\
U_{\mathrm{pred}},\
U_{\mathrm{stability}}).
\]

Интерпретация:

\begin{itemize}
    \item $U_{\mathrm{att}}$ — ландшафт потенциала, формирующий циклы внимания,
    \item $U_{\mathrm{wm}}$ — энергия, необходимая для поддержания состояний рабочей памяти,
    \item $U_{\mathrm{concept}}$ — концептуальный потенциал, разделяющий аттракторы,
    \item $U_{\mathrm{pred}}$ — потенциал предикативных рассогласований,
    \item $U_{\mathrm{stability}}$ — глобальный стабилизирующий потенциал геометрии аттракторов.
\end{itemize}

Когнитивный континуум требует:

\[
U_{\mathrm{concept}} \text{ имеет несколько минимумов, соответствующих понятиям}.
\]

% ---------------------------------------------------------------
\subsubsection*{4. Пороги \texorpdfstring{$\Theta(M_6)$}{\Theta(M_6)}}

Ключевые пороги включают:

\begin{itemize}
    \item $\Theta_{\mathrm{att}}$ — минимальная активация для вовлечения внимания,
    \item $\Theta_{\mathrm{wm}}$ — порог устойчивого поддержания памяти,
    \item $\Theta_{\mathrm{sep}}$ — порог разделимости между аттракторами,
    \item $\Theta_{\mathrm{concept}}$ — порог возникновения отдельных понятий,
    \item $\Theta_{\mathrm{coh}}$ — порог когнитивной когерентности.
\end{itemize}

Критическое условие существования когнитивного континуума:

\[
\Theta_{\mathrm{sep}} < \Delta_{\mathrm{basin}},
\quad
\Theta_{\mathrm{coh}} \le T_{\mathrm{cog}},
\]
где $\Delta_{\mathrm{basin}}$ — расстояние между аттракторами.

% ---------------------------------------------------------------
\subsubsection*{5. Потоки \texorpdfstring{$J(M_6)$}{J(M_6)}}

Допустимые когнитивные потоки включают:

\[
J(M_6) =
\big(
J_{\mathrm{att}},\
J_{\mathrm{wm}},\
J_{\mathrm{concept}},\
J_{\mathrm{pred}},\
J_{\mathrm{plasticity}},\
\nabla_i T_6
\big).
\]

Интерпретация:

\begin{itemize}
    \item $J_{\mathrm{att}}$ — поток модуляции внимания,
    \item $J_{\mathrm{wm}}$ — поток стабилизации рабочей памяти,
    \item $J_{\mathrm{concept}}$ — потоки, трансформирующие геометрию аттракторов,
    \item $J_{\mathrm{pred}}$ — распространение ошибки прогноза,
    \item $J_{\mathrm{plasticity}}$ — обновление синаптических весов, поддерживающее когерентность.
\end{itemize}

Соотношение сохранения:

\[
\sum_i J_{\mathrm{plasticity},i} = \frac{d\mathcal{W}}{dt}.
\]

% ---------------------------------------------------------------
\subsubsection*{6. Циклы \texorpdfstring{$C(M_6)$}{C(M_6)}}

Ключевые циклы, обеспечивающие когнитивную деятельность:

\[
C_{\mathrm{att\_wm}}
=
\{\text{внимание} \to \text{кодирование} \to
\text{поддержание} \to \text{обновление} \to \text{освобождение}\}.
\]

Дополнительные циклы:

\begin{itemize}
    \item $C_{\mathrm{predictive}}$ — предсказание $\to$ ошибка $\to$ обновление,
    \item $C_{\mathrm{concept}}$ — стабилизация концептуальных границ,
    \item $C_{\mathrm{proto\_logic}}$ — логические преобразования состояний,
    \item $C_{\mathrm{plasticity}}$ — циклы обучения, корректирующие $\mathcal{W}$.
\end{itemize}

Необходимое условие когнитивной жизни:

\[
\oint_{C_{\mathrm{att\_wm}}} dA_{\mathrm{state}} \approx 0.
\]

% ---------------------------------------------------------------
\subsubsection*{7. Граница \texorpdfstring{$\partial\Omega(M_6)$}{\partial\Omega(M_6)}}

Конфигурация достигает $\partial\Omega(M_6)$, когда:

\begin{itemize}
    \item аттракторы теряют устойчивость или коллапсируют,
    \item семантические оси становятся вырожденными,
    \item ошибка прогноза расходится (неограниченно),
    \item пластичность разрушает когерентность,
    \item $\Theta_{\mathrm{concept}}$ невозможно пересечь (нет формирования понятий),
    \item внимание не может стабилизироваться (цикл не срабатывает).
\end{itemize}

Пересечение границы означает коллапс $K_6$ в динамику типа $K_5$ (чисто электрическую).

% ---------------------------------------------------------------
\subsubsection*{8. Связь с \texorpdfstring{$M_5$}{M_5} и \texorpdfstring{$M_7$}{M_7}}

\textbf{От $M_5$ к $M_6$:}

Требования:

\[
C_{\mathrm{spike}} \text{ должно связываться со структурированными паттернами},
\quad
\mathcal{W} \text{ должна поддерживать бассейны притяжения}.
\]

Так последовательности спайков приобретают геометрию с смысловым содержанием.

\textbf{В сторону $M_7$ (социальная когниция):}

\[
A_{\mathrm{shared}} \in A(M_7),
\qquad
\text{требуется устойчивый концептуальный репертуар в } M_6,
\]
образуя основу коммуникации, общего внимания и социального вывода.
